%% kapa: exam=1 round=23 year=2012 period=2 form=A subject=appraisal_law subject_ko=부동산관계법규 seq=34 number=74
%% source: Q-Net 감정평가사 기출문제 (cst00309) artlSeq=5052580 file=raw/1008/5052580/2012년 제23회 감정평가사 제1차시험 기출문제(2교시 A형).pdf
%% conversion: kordoc 4.24.51 -> tools/gen1.py (auto); review=proofread-sonnet; flags=none
\begin{question}{74}
측량ㆍ수로조사 및 지적에 관한 법령상 지적측량을 하여야 하는 경우에 해당하지 않는 것은? (단, 측량의 필요성은 인정되고 도시개발사업 등의 시행지역이 아님)
\begin{choices}
  \item 지적공부를 복구하는 경우
  \item 토지를 등록전환하는 경우
  \item 축척을 변경하는 경우
  \item 토지를 합병하는 경우
  \item 경계점을 지상에 복원하는 경우
\end{choices}
\end{question}
\answer{4}
